\begin{table}[H]
\centering
\caption{Phân loại khoảng trống nghiên cứu và phạm vi xử lý của khoá luận}
\label{tab:v2-gap}
\renewcommand{\arraystretch}{1.25}
\small
\begin{tabular}{|c|p{2.2cm}|p{5.6cm}|p{4.6cm}|}
\hline
\textbf{Mã} & \textbf{Tầng} & \textbf{Khoảng trống} & \textbf{Khoá luận} \\ \hline
G1 & Chuẩn hoá & Chuẩn truy vấn sổ đăng ký hoãn phần uỷ quyền nhiều cấp \cite{ToIP_TRQP2026}; chuẩn trạng thái không biểu diễn thẩm quyền & Góp một cách biểu diễn cụ thể chạy được \\ \hline
G2 & Kiến trúc & Sổ đăng ký là cơ sở dữ liệu của một tổ chức, chỉ di chuyển sự phụ thuộc \cite{SSI_Survey2025} & Xử lý trực tiếp: sổ đăng ký là trạng thái hợp đồng \\ \hline
G3 & Biểu diễn thẩm quyền & Thẩm quyền nhị phân, thiếu phạm vi, thời hạn, số cấp & Xử lý trực tiếp: mặt nạ bit, thời hạn, uỷ quyền hai cấp \\ \hline
G4 & Ngữ nghĩa xác minh & Không phân biệt loại thất bại; không đối chiếu nội dung với băm đã neo & Xử lý trực tiếp: bảy bước báo cáo riêng \\ \hline
G5 & Phương pháp đánh giá & Chỉ 2/47 công trình có nghiên cứu người dùng \cite{Ambati2026} & Không xử lý, cùng chung giới hạn \\ \hline
G6 & Bối cảnh & Pháp lý và hạ tầng đã có, thiếu lớp quản trị thẩm quyền cho giáo dục & Xử lý một phần, trong một trường hợp nghiên cứu \\ \hline
\end{tabular}
\end{table}
